\documentclass{beamer}

\usepackage{amsmath,amssymb,amsthm}
\usepackage[all]{xy}
\usepackage{graphicx}
\usepackage[backend=biber,style=alphabetic,maxnames=99,giveninits=true,url=false,isbn=false]{biblatex}
\addbibresource{refs.bib}



\numberwithin{equation}{subsection}

% Basic macros
\def\lim{\mathop{\mathrm{lim}}\nolimits}
\def\colim{\mathop{\mathrm{colim}}\nolimits}
\def\Spec{\mathop{\mathrm{Spec}}}
\def\Hom{\mathop{\mathrm{Hom}}\nolimits}
\def\Ext{\mathop{\mathrm{Ext}}\nolimits}
\def\Ker{\mathop{\mathrm{Ker}}}
\def\Im{\mathop{\mathrm{Im}}}
\def\Coker{\mathop{\mathrm{Coker}}}
 \def\Coim{\mathop{\mathrm{Coim}}}

% Slide numbering
\addtobeamertemplate{navigation symbols}{}{%
    \usebeamerfont{footline}%
    \usebeamercolor[fg]{footline}%
    \hspace{1em}%
    \insertframenumber/\inserttotalframenumber
}



\title{Deformations of Stanley-Reisner schemes}
\author{Daniel González Casanova Azuela}
\institute{PUC-Rio}
\date{Informal talk\\ PUC-Rio\\[1em] December 3, 2026}

\begin{document}

\begin{frame}%Title page
\titlepage
\end{frame}

\section*{Question}

\begin{frame}%Question
\frametitle{The question}
Let $\mathcal{M}$ be Grünbaum's triangulated $3$-sphere with $8$ vertices.

\bigskip\noindent
\textbf{Does the Stanley-Reisner scheme of $\mathcal{M}$ deform
to a smooth Calabi-Yau threefold of degree $20$ in $\mathbb{P}^7$?}

\bigskip\noindent
\textbf{Today:} a proposed proof that the answer is \emph{yes},
and one local case of it done completely.

\bigskip\noindent
{\color{red}Status: the proof was assembled with AI help and
checked only by an AI audit. It is not yet verified by a human referee.
Red = not verified by me.}
\end{frame}

\begin{frame}%Plan
\frametitle{Plan}
\tableofcontents
\end{frame}

\section{Stanley-Reisner schemes}

\begin{frame}%SR schemes
\frametitle{Stanley-Reisner schemes}
A \emph{simplicial complex} $\mathcal{K}$ on $\{0,\ldots,n\}$:
a set of subsets (\emph{faces}) closed under taking subsets.

\medskip\noindent
Its ideal is generated by the monomials of the \emph{non-faces}:
\begin{align*}
I_{\mathcal{K}}&=\left<x_{i_1}\cdots x_{i_k} : \{i_1,\ldots,i_k\}\notin\mathcal{K}\right>
\subseteq \mathbb{C}[x_0,\ldots,x_n],\\
\text{SR}(\mathcal{K})&=\text{Proj}\,\mathbb{C}[x_0,\ldots,x_n]/I_{\mathcal{K}}.
\end{align*}

\medskip\noindent
\textbf{Example.} The path $\bullet-\bullet-\bullet-\bullet$ with edges
$01,12,23$. Non-faces: $02,03,13$, so
$$
I=\langle x_0x_2,\ x_0x_3,\ x_1x_3\rangle .
$$
$\text{SR}(\mathcal{K})$ is the union of three lines in $\mathbb{P}^3$,
one for each edge.

\medskip\noindent
In general, $\text{SR}(\mathcal{K})$ is a union of coordinate subspaces,
one for each maximal face.
\end{frame}

\section{Smoothings}

\begin{frame}%Smoothing
\frametitle{Smoothings}
A \emph{deformation} of $X$ is a flat family $\mathcal{X}\to B$
with $X$ as the fibre over a point $0\in B$.
A \emph{smoothing} is a deformation whose general fibre is smooth.

\bigskip\noindent
\textbf{Basic example.}
$$
\mathcal{X}=\{xy=t^2\}\subset\mathbb{A}^2\times\mathbb{A}^1_t .
$$
\begin{itemize}
\item For $t\neq0$ the fibre $xy=t^2$ is a smooth conic.
\item For $t=0$ the fibre $xy=0$ is two lines meeting at the origin.
\end{itemize}

\medskip\noindent
Why every fibre with $t\neq0$ is smooth: a singular point would have
$x=y=0$, which gives $0=t^2$.
\end{frame}

\begin{frame}%Twisted cubic
\frametitle{The path smooths to the twisted cubic}
The twisted cubic has ideal
$$
x_1^2-x_0x_2,\quad x_1x_2-x_0x_3,\quad x_2^2-x_1x_3 .
$$
Substitute $x_i\rightsquigarrow t^{w_i}x_i$ with $w=(1,2,2,0)$:
$$
t^{4}x_1^2-t^{3}x_0x_2,\qquad t^{4}x_1x_2-t^{1}x_0x_3,\qquad t^{4}x_2^2-t^{2}x_1x_3 .
$$
Divide each generator by its lowest power of $t$. At $t=0$
only the second terms survive:
$$
x_0x_2,\qquad x_0x_3,\qquad x_1x_3 .
$$
So the three lines are a flat limit of the twisted cubic:
the twisted cubic is a smoothing of the path.
The flatness is the theory of Gröbner degenerations.
\end{frame}

\section{The Grünbaum sphere}

\begin{frame}%Why spheres
\frametitle{Why spheres}
If $\mathcal{K}$ triangulates a $3$-sphere, $\text{SR}(\mathcal{K})$ is
a union of $\mathbb{P}^3$'s glued along that sphere,
of degree the number of facets.

\medskip\noindent
Altmann--Christophersen \cite{AC} relate the smoothings of such
schemes to Calabi-Yau geometry.

\medskip\noindent
{\color{red}Precise statement: not rechecked by me.
Today's theorem proves the Calabi-Yau properties directly,
for one complex.}
\end{frame}

\begin{frame}%Grünbaum sphere
\frametitle{Grünbaum's sphere}
Grünbaum--Sreedharan \cite{GS}: there is a triangulation $\mathcal{M}$ of
the $3$-sphere with $8$ vertices that is \emph{not} the boundary of a
$4$-polytope.

\medskip\noindent
Vertices $a,\ldots,h$. Its Stanley-Reisner ideal is generated by sixteen cubics:
$$
\begin{aligned}
&abf,\ abg,\ abh,\ acg,\ ach,\ adh,\ bdf,\ bdg,\\
&beg,\ cde,\ ceg,\ ceh,\ cfh,\ def,\ dfh,\ efg .
\end{aligned}
$$
$\text{SR}(\mathcal{M})\subset\mathbb{P}^7$ has dimension $3$ and degree $20$.

\medskip\noindent
Symmetry: $G\simeq S_3$ generated by $(b\,h)(c\,g)(d\,f)$ and
$(a\,c)(b\,e)(d\,f)$. Vertex orbits: $\{a,c,g\}$, $\{b,e,h\}$, $\{d,f\}$.
\end{frame}

\section{The result}

\begin{frame}%Main theorem
\frametitle{The result {\color{red}(proposed proof)}}
\begin{theorem}[{\color{red}proposed}]
\begin{enumerate}
\item The ideal $I_{\mathcal{M}}$ has an $S_3$-equivariant flat deformation
over $\mathbb{C}[[t]]$ whose generic fibre is smooth.
\item It gives an algebraic smoothing of $\text{SR}(\mathcal{M})$ over a curve.
\item The smooth fibres $Y\subset\mathbb{P}^7$ satisfy
\begin{gather*}
K_Y\simeq\mathcal{O}_Y,\quad H^1(\mathcal{O}_Y)=H^2(\mathcal{O}_Y)=0,\\
H^3=20,\quad c_2\cdot H=56.
\end{gather*}
\end{enumerate}
\end{theorem}
\noindent
Not claimed: Hodge numbers, fundamental group, uniqueness, or
that $Y$ is the Gulliksen-Negård threefold.
\end{frame}

\begin{frame}%Strategy
\frametitle{Strategy}
\begin{enumerate}
\item \textbf{Deformation theory.} The obstruction space has no $S_3$-invariants,
so every equivariant first-order deformation extends to all orders.
\item \textbf{A first-order deformation.} Ten monomial maps, summed
over $S_3$-orbits.
\item \textbf{The second-order terms are forced} by flatness.
\item \textbf{No singular points.} A singular point of the generic fibre would
specialize to a point of $\text{SR}(\mathcal{M})$. Rule this out
face by face: triangles, edges, vertices.
\item \textbf{Curve selection} turns the formal family into an algebraic one;
Riemann-Roch gives the invariants.
\end{enumerate}
\end{frame}

\section{Local case: a triangle}

\begin{frame}%Step 1
\frametitle{Step 2: the first-order deformation}
A first-order deformation replaces each generator $m_i$ by $m_i+tH_i$.
Each of the ten maps has the form
$$
\phi_{p,q,J}(m_i)=\frac{p\,m_i}{q}\quad(i\in J),\qquad 0\quad(i\notin J).
$$

\medskip\noindent
\textbf{Example.} $p=ce$, $q=bf$, $J=\{1\}$, $m_1=abf$:
$$
abf\ \longmapsto\ \frac{ce\cdot abf}{bf}=ace .
$$
Flatness needs: for $j\notin J$, $p\cdot\text{lcm}(m_1,m_j)/q\in I$.
Here it comes down to
\begin{align*}
ace\cdot d&\in(cde),&ace\cdot g&\in(ceg),&ace\cdot h&\in(ceh).
\end{align*}
\end{frame}

\begin{frame}%Local equations
\frametitle{The open triangle $beh$}
Near the points with $b,e,h\neq0$ and all other coordinates zero, set
$b=1$, $e=u$, $h=v$ (units), $d=U$, $f=V$.

\medskip\noindent
The special fibre becomes $(a,c,g,\,UV)$: the cubics
$abh,\ ceh,\ beg,\ bdf$ become $va,\ uvc,\ ug,\ UV$.

\medskip\noindent
The first three equations are solved for $a,c,g$ (their
derivative matrix is $\text{diag}(v,uv,u)$, a unit). What remains is
$$
f(U,V,u,v,t)\;=\;UV\;+\;t^2\,p(u,v)\;+\;\cdots
$$
with \emph{no $t$-term}, and
$$
p(u,v)=-8u^2-20uv-8v^2-20u-20v-8.
$$
{\color{red}(From the manuscript; $p$ is forced by flatness.)}
\end{frame}

\begin{frame}%Triangle argument
\frametitle{No singular point near the triangle}
The model is $xy=t^2$, with $UV+t^2p$ in place of $xy-t^2$.

\medskip\noindent
A singular point has critical value $\Delta=t^2p(u,v)+\cdots$ and
$\Delta_u=\Delta_v=0$. Dividing by $t^2$ and setting $t=0$ gives
a common zero of $p,\,p_u,\,p_v$.

\medskip\noindent
But
$$
p-\tfrac12\bigl(u+\tfrac59\bigr)p_u-\tfrac12\bigl(v+\tfrac59\bigr)p_v=\tfrac{28}{9},
$$
which is a unit in characteristic $0$ and modulo $101$.
So $p,p_u,p_v$ have no common zero: no singular point.

\medskip\noindent
{\color{red}Not verified by me: the derivation of $f$ and of the critical-value argument.}
\end{frame}

\section{What remains}

\begin{frame}%Rest of the proof
\frametitle{The other cases}
The same scheme applies to every face of $\mathcal{M}$.
\begin{itemize}
\item \textbf{Edges:} the case $be$ has two ranges of valuations and an
exceptional cancellation, $H_7=10acd+10cfg$.
\item \textbf{Other triangles and edges:} tables and short calculations.
\item \textbf{Vertices:} three orbits, $a,\ b,\ d$.
\end{itemize}
{\color{red}Everything above is verified only by AI audit.}
\end{frame}

\begin{frame}%Open
\frametitle{Open questions}
\begin{itemize}
\item Is the smoothing the Gulliksen-Negård threefold? Open: a comparison with a
del Pezzo construction suggested by Grzegorz is unresolved.
\item Hodge numbers, fundamental group?
\item Other non-polytopal spheres.
\item Human verification of the proof.
\end{itemize}

\medskip\noindent
\textbf{Questions for you:} (1) Which step would you check first?
(2) Do you know a conceptual reason why this complex smooths?
\end{frame}

\begin{frame}[allowframebreaks]%References
\frametitle{References}
\printbibliography[heading=none]
\end{frame}

\end{document}
